\multiproblem{discrete}{Let us consider some simple questions on discrete random variables to warm up.
\begin{enumerate}
 \item Consider the random trial of tossing a fair die. Let $X$ be the number on the die.
 \begin{enumerate}
  \item Define and sketch the probability function and cdf.
  \item Find the expected value $E(X)$ and variance $\text{Var}(X)$.
 \end{enumerate}
\item Now imagine that our die is no longer fair. Instead it has the pdf given by the following table
\begin{center}
\begin{tabular}{ |c|c|c|c|c|c|c| } 
 \hline
 $x$ & 1 & 2 & 3& 4& 5& 6  \\ 
\hline $P_{X}(x)$ & $a$ & $a$ & $a$ & $b$ & $b$ & 0.3   \\ 
 \hline
\end{tabular}
\end{center}
\begin{enumerate}
 \item Find the values of $a$ and $b$, given that $E(X)=4.2$. 
 \item Sketch the pdf and cdf.
 \item Find $\text{Var}(X)$.
\end{enumerate}
\end{enumerate}
}

\multiproblem{cts}{Now let us consider some simple examples of continuous random variables.
\begin{enumerate}
\item A continuous \emph{uniform} random variable $X$ has the following pdf
\begin{equation*} f_X(x) = \begin{cases} 
      \frac{1}{a}  &  0 \leq x \leq a \\
      0 &  \text{else} 
        		\end{cases} \end{equation*}
        		where $a$ is a constant scalar.
        		\begin{enumerate}
        		 \item Sketch its pdf and its cdf.
        		 \item Find $E(X)$ and $\text{Var}(X)$.
        		\end{enumerate}

\item The pdfs of two different continuous random variables are shown below.\\
\begin{center}
        \includegraphics[scale=2.5]{pdfs.pdf}
        \end{center}
        For each:
        \begin{enumerate}
         \item Find the value of the constant scalar $a$.
         \item Find the pdf.
         \item Find $E(X)$ and $\text{Var}(X)$.
        \end{enumerate}
\end{enumerate}
}
   
\multiproblem{proofs}{Let us now consider some more proofy questions. Consider random variable $X$ and constant scalar $a$.
\begin{enumerate}
 \item When $X$ is continuous, prove that $E(aX)=aE(X)$. Do the same for when $X$ is discrete.
 \item Prove that $\text{Var}(a X)=a^{2}\text{Var}(X)$.
\end{enumerate}
}

\multiproblem{binom1}{
\begin{enumerate}
\item A fair die is thrown 12 times. Find the probability that the die will show a number less than 5 on exactly 4 of the throws.
\item One week Alice received 25 spam emails and 75 non-spam emails. She then received 10 emails. 
\begin{enumerate}
 \item Estimate the probability that exactly 6 of those will be spam?
 \item Estimate the probability that she will receive less than 3 spam emails?
\end{enumerate}
\end{enumerate}
}

\multiproblem{binom2}{
A Bernoulli random variable has two possible outcomes: 0 or 1. It is
characterised by the probability $p$ of having the outcome 1. For example
the random variable representing the number of heads from flipping a fair coin
once is a Bernoulli random variable with $p=\frac{1}{2}$.

A binomial distributed random variable is characterised by two parameters $N$
and $p$. It represents the sum of $N$ independent and identically distributed
Bernoulli random variables each having probability $p$. For example the number
of heads obtained from flipping a fair coin 3 times is binomially distributed
with $N=3$ and $p=\frac{1}{2}$.
\begin{enumerate}
\item The probability that a certain kind of component will survive a shock
    test is 3/4. Use the binomial distribution to estimate the probability
    that exactly 2 of the next 4 components tested survive.
\item The probability that a patient recovers from a rare blood disease is
    0.4. 15 people are known to have contracted this disease. Assuming that
    their survival probabilities are independent, what is the probability that
\begin{enumerate}
 \item at least 10 survive
 \item from 3 to 8 survive
 \item exactly 5 survive?
\end{enumerate}
When might the assumption of independent survival probabilities not be
satisfied?
\item It is conjectured that an impurity exists in 30\% of all drinking wells in a certain rural community. In order to gain some insight into the true extent of the problem, it is determined that some testing is necessary. It is too expensive to test all of the wells in the area, so 10 are randomly selected for testing.
\begin{enumerate}
 \item Using the binomial distribution, estimate is the probability that exactly 3 wells have the impurity, assuming that the conjecture is correct?
 \item What would be the probability that more than 3 wells are impure?
\end{enumerate}
\end{enumerate}
}

\multiproblem{expon}{
\begin{enumerate}
\item Prove that the expected value of an exponential random variable $X$ is $E(X)=\frac{1}{\lambda}$
\item Prove that the variance of an exponential random variable is $\text{Var}(X)=\frac{1}{\lambda^2}$
\item Let $X$ be an exponential random variable with parameter $\lambda=\ln(3)$. Compute the following probability: $P(3 \leq X \leq 4)$.
\item Suppose that a system contains a certain type of component whose time, in years, to failure is given by T. The random variable T is modeled nicely by the exponential distribution with mean time to failure of 5 years. If 5 of these components are installed in different systems, what is the probability that at least 2 are still functioning at the end of 8 years?
\item The life of a disposable pen satisfies an exponential distribution with a mean and standard deviation of 30 hours. A group of 50 students are each given a new pen at the start of a 2 hour maths exam.
\begin{enumerate}
 \item Find the probability that a typical student will need at least one replacement pen to finish the exam.
 \item What is the probability that four or more students need a replacement pen?
\end{enumerate}
\end{enumerate}
}

\multiproblem{poisson}{
\begin{enumerate}
\item Prove that the expected value of the Poisson distribution is $\lambda$.
\item During a laboratory experiment, the average number of radioactive particles passing through a counter in 1 millisecond is 4. What is the probability that 6 particles enter the counter in a given millisecond?
\item Ten is the average number of oil tankers arriving each day at a certain
    port. The facilities at the port can handle at most 15 tankers per day.
    Estimate the probability that on a given day tankers have to be turned
    away.\\

If n is quite large and p is small we can use the Poisson distribution to
approximate the Binomial distribution.

 \item In a certain industrial facility, accidents occur infrequently. It is
     known that the probability of an accident on any given day is 0.005 and
     accidents are independent of each other.
\begin{enumerate}
 \item What is the probability that in any given period of 400 days there will
     be an accident on exactly one day?
 \item What is the probability that there are at most three days with an accident?
\end{enumerate}
For both of the above, first calculate the probability using the Binomial
distribution. Then find the Poisson approximation and compare the two values.
\end{enumerate}
}

\multiproblem{normal}{
\begin{enumerate}
\item The tensile strength of a certain metal component is normally distributed with a mean of 10,000 kilograms per square centimeter and a standard deviation of 100 kilograms per square centimeter. Measurements are recorded to the nearest 50 kilograms per square centimeter.
\begin{enumerate}
 \item What proportion of these components exceed 10,150 kilograms per square centimeter in tensile strength?
 \item If specifications require that all components have tensile strength between 9800 and 10,200 kilograms per square centimeter inclusive, what proportion of pieces would we expect to scrap?
\end{enumerate}
 \item  A certain type of storage battery lasts, on average, 3.0 years with a standard deviation of 0.5 year. Assuming that battery life is normally distributed, find the probability that a given battery will last less than 2.3 years.
\end{enumerate}
}
